\documentclass[../../main.tex]{subfiles}
\begin{document}

\begin{center}
  {\fontsize{48pt}{58pt}\selectfont \textbf{1973}\normalsize\textbf{年度}}
\end{center}

\begin{table}[htb]
  \centering
  \caption{出題テーマと難易度}
  \begin{tabular}{|c|p{9cm}|c|}
    \hline
    \textbf{問題番号} & \textbf{テーマ} & \textbf{難易度} \\
    \hline
    第1問 & 整数問題（不等式） & \\
    \hline
    第2問 & 整式の剰余と決定 & \\
    \hline
    第3問 & 三角関数の不等式証明 & \\
    \hline
    第4問 & 三次関数の極値 & \\
    \hline
    第5問 & 回転体の体積の一定性証明 & \\
    \hline
    第6問 & 確率 & \\
    \hline
  \end{tabular}
\end{table}

\begin{table}[htb]
  \centering
  \caption{解答}
  \renewcommand{\arraystretch}{1.4}
  \begin{tabular}{|c|c|p{9cm}|}
    \hline
    \textbf{問題番号} & \textbf{小問} & \textbf{答え} \\
    \hline
    第1問 & - & $(p,q)=(0,\pm1),(\pm1,\pm1),(\pm1,\pm2)$（複号同順） \\
    \hline
    第2問 & - & $f(x)=-2x^3-3x^2+x+1$ \\
    \hline
    第3問 & - & 証明問題（$x=1$を示す） \\
    \hline
    第4問 & - & $p=-\dfrac{a^2}{3},\ q=\dfrac{2}{27}a^3+b$ \\
    \hline
    第5問 & - & 証明問題（$ab=\dfrac{5}{6}$） \\
    \hline
    第6問 & - & $\dfrac{17}{108}$ \\
    \hline
  \end{tabular}
\end{table}

\paragraph{試験形態}
6問．試験時間は要確認．

\end{document}
